\begin{abstract}
方向消融能够在不重新训练基础模型的前提下改变大语言模型的拒答行为，但人工选择方向、作用层和消融强度既费时，也容易牺牲正常行为。本文基于固定源码提交，对一套自动化方向消融流程进行可追踪分析：系统从有害与无害提示的首次响应残差中提取候选拒答方向，把注意力输出投影与多层感知机下投影的作用范围参数化，并用关键词拒答率和首响应词元分布差异构成多目标搜索，最后导出低秩适配器或合并权重。针对 Qwen3.8-27B，本文说明其门控线性注意力与门控注意力如何被统一映射到消融组件，并审查两份公开实践记录。主记录的第三方模型卡报告关键词拒答由 99/100 降至 4/100，首响应词元分布差异为 0.0796；这些数字不是本文独立实验。分析表明，该流程提供了从源码机制到公开案例的审计链，但其结论受到关键词代理、单词元测量、四比特搜索与脑浮点十六位权重合并差异、思考模板切点以及未评估视觉路径的共同限制。本文据此把该系统定位为可解释的模型行为干预工具，而非能力无损或安全合规的证明。
\end{abstract}

\begin{IEEEkeywords}
大语言模型，方向消融，拒答行为，多目标优化，Qwen3.8
\end{IEEEkeywords}
